\documentclass[11pt]{article}

\usepackage[a4paper,margin=30mm]{geometry}
\usepackage{amsmath,amssymb,amsthm,mathtools}
\usepackage{booktabs,tabularx,array,float}
\usepackage{enumitem}
\usepackage{microtype}
\usepackage[hidelinks]{hyperref}

\newtheorem{theorem}{Theorem}[section]
\newtheorem{proposition}[theorem]{Proposition}
\newtheorem{corollary}[theorem]{Corollary}
\newtheorem{lemma}[theorem]{Lemma}
\theoremstyle{definition}
\newtheorem{definition}[theorem]{Definition}
\newtheorem{problem}[theorem]{Open source law}
\theoremstyle{remark}
\newtheorem{remark}[theorem]{Remark}

\newcommand{\C}{\mathbb C}
\newcommand{\R}{\mathbb R}
\newcommand{\N}{\mathbb N}
\newcommand{\Z}{\mathbb Z}
\newcommand{\XiC}{\mathcal X}
\newcommand{\RH}{\mathrm{RH}}
\newcommand{\Div}{\operatorname{Div}}
\newcommand{\ord}{\operatorname{ord}}
\newcommand{\ran}{\operatorname{ran}}
\newcommand{\mul}{\operatorname{mul}}
\newcommand{\dom}{\operatorname{dom}}
\newcommand{\Span}{\operatorname{span}}
\newcommand{\Rea}{\operatorname{Re}}
\newcommand{\Ima}{\operatorname{Im}}

\title{\textbf{Arithmetic source descent toward the Riemann hypothesis}\\
\large Cyclic residue transgression and maximal quotient collapse}
\author{Michael Haufschild}
\date{}

\begin{document}
\maketitle

\begin{abstract}
We analyze a source-first route to the Riemann hypothesis and separate its
constructed transport from its selection problem.  A target-free
M\"obius--theta source reaches every jet of the centered completed zeta
function.  At a zero of arbitrary multiplicity it retains the first nonzero
jet as a division-free top conormal, lands coefficient-faithfully in a
completed Gamma/co-Poisson incidence, and survives a boundedly invertible
reflected Green transport.  Its unsquared Hodge trace has exact radical
\(\Rea\rho=1/2\).  None of these facts excludes an off-line divisor.

We prove a universal boundary-defect law.  For a cyclic source rotation
\(\tau\), paracyclic target rotation \(T\), landing \(L\), and defect
\(D=TL-L\tau\),
\[
                         TD+D\tau=T^2L-L\tau^2.
\]
If \(\tau^2=1\) and \(T^2=\delta\,\mathrm{id}\), strict equivariance forces
every off-unit-monodromy landing to vanish.  The smallest admissible boundary
is
\[
                         B_{\min}=\ran D+\ran(TD).
\]
It is target-stable, contains every defect, and lies in every competing
stable defect boundary; hence it gives the largest possible quotient.  Away
from \(\delta=1\), the landing range lies in \(B_{\min}\).

We then construct the previously missing comparison from the literal
degree-one rational-adelic cyclic bar to the physical two-leg landing,
before choosing a zero.  For every ordinary prime \(p\), a canonical
\(p\otimes p^{-1}\) cycle has trace one.  The coordinate-sum cyclic residue
kills its untwisted cyclic comparison defect, detects its physical leading
jet, and retains multiplicity in a separate jet filtration.  Under the
analytic paracyclic rotation its defect satisfies the exact transgression
law
\[
 \operatorname{Res}(D_{n,\rho}x)
  =(\delta_n(\rho)-1)c_\rho\,\varepsilon_0^{(1)}(x).
\]
For the physical two-leg landing we prove
\(B_{\min,\rho}=0\) on the critical line and
\(B_{\min,\rho}=\C^2\) off it.  Consequently the constructed residue
annihilates \(B_{\min,\rho}\) exactly on the critical line, and global
annihilation is equivalent to RH.  Thus source comparison, ordinary-prime
detection, multiplicity retention, boundary choice, and singular topology
are closed unconditionally.  The sole RH-bearing content is now the
source-null-homology of the physical paracyclic defect; this paper neither
assumes nor proves that final statement.
\end{abstract}

\medskip
\noindent\textbf{Keywords.}
Riemann hypothesis; M\"obius--theta source; conormal; Gamma/co-Poisson
boundary; source descent; cyclic monodromy; singular completion.

\section{The problem is selection, not representation}

Write
\[
 \xi(s)=\frac12s(s-1)\pi^{-s/2}\Gamma(s/2)\zeta(s),
 \qquad
 \XiC(w)=\xi\!\left(\frac12+iw\right).
\]
For a nontrivial zero \(\rho\) of \(\xi\), put
\[
 w_\rho=(\rho-\tfrac12)(-i),
 \qquad
 a_\rho=\Rea\rho-\frac12.
\]
Then \(\XiC(w_\rho)=0\), and the Riemann hypothesis is the assertion
\(a_\rho=0\) for every divisor point.  Let \(r_\rho\geq1\) be the analytic
multiplicity of \(\XiC\) at \(w_\rho\).  Thus
\[
 \XiC^{(j)}(w_\rho)=0\quad(0\leq j<r_\rho),
 \qquad
 \XiC^{(r_\rho)}(w_\rho)\neq0.
\tag{1.1}\label{eq:first-jet}
\]

The constructions below show that the arithmetic source reaches the entire
local Cartier package \eqref{eq:first-jet}, retains its leading conormal, and
places it in the Gamma/co-Poisson incidence geometry without annihilation.
Consequently the unresolved issue is not whether a source reaches a zero.
It is why one particular real character carried by that nonzero source atom
must be unitary.

The distinction matters because a representation theorem and a selection
theorem have opposite logical directions.  Representation constructs a
carrier for every possible divisor, including a hypothetical off-line one.
Selection proves that the arithmetic source cannot support the off-line
character.  A faithful representation can therefore be complete and still
make no progress on zero exclusion.

\section{The source-descent conservation law}

We first state the universal mechanism used to audit every transport in the
route.

\begin{definition}
Let \(T:E\to F\) be a linear map and \(S\subseteq E\) a source sector.  We
say that \(T\) is \emph{conservative on \(S\)} when
\[
                 S\cap\ker T=\{0\}.
\]
For a grade map \(g:I\to\Gamma\) and a finite packet
\((c_i)_{i\in A}\), its resolved realization is
\[
 P(\gamma)=\sum_{\substack{i\in A\\g(i)=\gamma}}c_i.
\]
\end{definition}

\begin{theorem}[Source-descent conservation]
\label{thm:source-descent}
Let \(E\xrightarrow{T}F\xrightarrow{U}G\) be linear maps and
\(S\subseteq E\) a source sector.
\begin{enumerate}[label=\textup{(\roman*)}]
\item The composite \(UT\) is conservative on \(S\) if and only if
      \(T\) is conservative on \(S\) and \(U\) is conservative on
      \(T(S)\).
\item A finite resolved packet vanishes if and only if it vanishes in each
      full grade separately.  In particular, every nonzero term of a
      vanishing packet has a distinct, oppositely signed mate of the same
      full grade.
\item If a coarsening of grades creates a zero that was absent before
      coarsening, then two distinct full grades have collided.
\end{enumerate}
\end{theorem}

\begin{proof}
For (i),
\[
 \ker(UT)=T^{-1}(\ker U).
\]
Intersecting with \(S\) gives the forward implication.  Conversely, if
\(x\in S\) and \(UTx=0\), conservativity of \(U\) on \(T(S)\) gives
\(Tx=0\), and conservativity of \(T\) on \(S\) gives \(x=0\).

For (ii), equality in the product \(\prod_{\gamma\in\Gamma}A_\gamma\) is
coordinatewise.  If the coefficient at \(i\) is nonzero while its grade
coordinate sums to zero, at least one other nonzero coefficient occurs in
the same grade.  Part (iii) is the contrapositive: if the coarsening is
injective on the grades appearing in the packet, its fibre sums are exactly
the original grade coordinates and cannot create a new zero.
\end{proof}

\begin{corollary}[Forced order]
\label{cor:forced-order}
Suppose orientation is carried by a resolved source coordinate.  Any route
that first passes through an orientation-invariant quotient, a positive
quadratic response, or a completion whose graph has a nonzero vertical
fibre must recover orientation from an additional transverse channel.  No
later injective postprocessing of the reduced record repairs the lost edge.
\end{corollary}

This is not a Riemann-specific no-go.  It is the reason that the order
\[
 \boxed{\begin{gathered}
   \text{resolve source}\ \longrightarrow\
   \text{land analytically}\ \longrightarrow\
   \text{retain conormal}\\[-1mm]
   \longrightarrow\ \text{retain boundary class}\
   \longrightarrow\ \text{polarize}
 \end{gathered}}
\tag{2.1}\label{eq:forced-order}
\]
is forced rather than stylistic.

\section{A target-free arithmetic source reaches every Xi jet}

Let \(\Theta\) denote the fixed theta-cusp half-density on the logarithmic
line, and let \(\mu\) be the M\"obius function.  For an integer \(m\geq3\),
write \(\mu_m(n)\) for the component of \(\mu(n)\) supported on perfect
\(m\)-th powers.  Define the physical rung
\[
 \Psi_m(x)=\sum_{n\geq1}\mu_m(n)n^{-1/2}\Theta(x-\log n).
\tag{3.1}\label{eq:physical-rung}
\]
The series is absolutely convergent in the source Hilbert space.  It is
formed before a spectral point, a zero, or a divisor is chosen.

Internal polynomial promotion of \(\Theta\) gives sources
\(\Psi_{m,k}\) for every derivative order \(k\geq0\).  Translation on the
logarithmic line becomes multiplication by a Mellin character.  The M\"obius
Dirichlet series then gives the following exact transform.

\begin{proposition}[Physical source-to-Xi landing]
\label{prop:source-landing}
For \(|\Ima w|<1\) and
\(\Rea(\tfrac12+iw)>1/m\),
\[
 \mathcal{FL}(\Psi_{m,k})(w)
 =\zeta\!\left(m\left(\tfrac12+iw\right)\right)^{-1}
  \XiC^{(k)}(w).
\tag{3.2}\label{eq:source-transform}
\]
For every fixed \(w\) with
\(\Rea(\tfrac12+iw)>0\), the right-hand side converges to
\(\XiC^{(k)}(w)\) as \(m\to\infty\) through the integer rungs.
\end{proposition}

\begin{proof}
Absolute summability permits termwise Fourier--Laplace transform in
\eqref{eq:physical-rung}.  The logarithmic translate contributes
\(n^{-iw}\); the theta source contributes \(\XiC(w)\), and internal
promotion contributes its \(k\)-th derivative.  The remaining Dirichlet
series is
\[
 \sum_{n\geq1}\mu_m(n)n^{-1/2-iw}
 =\zeta\!\left(m\left(\tfrac12+iw\right)\right)^{-1},
\]
in its absolute-convergence domain, proving \eqref{eq:source-transform}.
For fixed positive real part, the scaled argument tends to the far right
half-plane and \(\zeta(ms)^{-1}\to1\).  The same domination gives convergence
of the transformed source.
\end{proof}

There is also a complete positive resolution.  For a base \(K\geq12\), let
\(M_{K,j}(s)\) be its high-exponent multiplier coordinates.  They satisfy
\[
 \sum_{j\geq0}|M_{K,j}(s)|<\infty,
 \qquad
 \sum_{j\geq0}M_{K,j}(s)=1
\tag{3.3}\label{eq:mass-one}
\]
throughout the closed right half of the critical strip.  Thus at least one
coordinate is nonzero at every point, and multiplication by that analytic
coordinate preserves exact analytic order.

\begin{theorem}[Cofinal conormal non-erasure]
\label{thm:conormal}
Let \(\rho\) be any nontrivial zero and \(r=r_\rho\).  The complete
M\"obius--theta source satisfies
\[
 \sum_{j=0}^{N}\mathcal{FL}(C_{K,j,k})(w_\rho)
 \longrightarrow \XiC^{(k)}(w_\rho)
\]
for every \(k\), where \(C_{K,j,k}\) are the exact adjacent complete-source
cofaces.  Its adjacent division-free conormal pair at the leading
order satisfies
\[
 \bigl(\mathcal{FL}(\Psi_{K+N,r-1})(w_\rho),
       \mathcal{FL}(\Psi_{K+N,r})(w_\rho)\bigr)
 \longrightarrow
 \bigl(0,\XiC^{(r)}(w_\rho)\bigr),
\tag{3.4}\label{eq:vertical-conormal}
\]
and the second coordinate is nonzero.
\end{theorem}

\begin{proof}
Apply Proposition~\ref{prop:source-landing} to the internally promoted
orders \(r-1\) and \(r\), use the mass-one telescoping
\eqref{eq:mass-one}, and pass to the high-exponent limit.  Equation
\eqref{eq:first-jet} gives the two limiting coordinates and proves
non-erasure.
\end{proof}

Theorem~\ref{thm:conormal} is stronger than pointwise value landing.  It
retains arbitrary multiplicity without dividing by \(\XiC\) or selecting a
nonzero coordinate after the fact.

\section{Coefficient-faithful Gamma/co-Poisson incidence}

Fix two distinct ordinary primes \(p\neq q\).  Let \(V_{p,q,\rho}\) be the
canonical completed Calder\'on--Witt incidence vector at the first
multiplicity copy of \(\rho\).  Normalize its distinguished coefficient to
one.  Define
\[
 c_\rho=\XiC^{(r_\rho)}(w_\rho),
 \qquad
 \mathcal L_{p,q}(\rho)=c_\rho V_{p,q,\rho}.
\tag{4.1}\label{eq:gamma-landing}
\]
The coefficient \(c_\rho\) is not an auxiliary divisor weight: it is the
diagonal physical moment of the internally promoted source.

\begin{theorem}[Global atomic Gamma incidence]
\label{thm:gamma-atomic}
For every nontrivial divisor point \(\rho\),
\[
 \mathcal L_{p,q}(\rho)\neq0.
\tag{4.2}\label{eq:gamma-nonzero}
\]
Decoding \(\mathcal L_{p,q}(\rho)\) gives
\[
 c_\rho\,W_{p,q},
\tag{4.3}\label{eq:decoded-witness}
\]
where \(W_{p,q}\) is the actual Gamma/Calder\'on--Witt incidence witness.
For every real modular displacement \(a\), the completed adjacent
Hellinger--Gamma cofaces recover \(\mathcal L_{p,q}(\rho)\) in the cofinal
depth limit.
\end{theorem}

\begin{proof}
The distinguished coefficient of \(V_{p,q,\rho}\) is one.  Therefore the
same coefficient of \(\mathcal L_{p,q}(\rho)\) is \(c_\rho\), which is
nonzero by \eqref{eq:first-jet}.  This proves \eqref{eq:gamma-nonzero}.
Linearity of the incidence decoder gives \eqref{eq:decoded-witness}.
The adjacent cofaces telescope to the normalized depth vertices; their
completed recovery map converges to the original incidence vector, and
scalar multiplication by \(c_\rho\) commutes with this limit.
\end{proof}

Let \(\mathfrak G_{p,q,n}(a)\) denote the unsquared secondary Hodge trace of
the gauged depth-\(n\) Gamma cell.  Its exact radical is
\[
                 \mathfrak G_{p,q,n}(a)=0
                 \quad\Longleftrightarrow\quad a=0.
\tag{4.4}\label{eq:hodge-radical}
\]

\begin{corollary}[Source-native unsquared detector]
\label{cor:unsquared-detector}
The Hodge trace of the physical landing is
\[
 \operatorname{Tr}_{H}\mathcal L_{p,q,n}(\rho,a)
 =c_\rho\,\mathfrak G_{p,q,n}(a),
\tag{4.5}\label{eq:joined-trace}
\]
and hence
\[
 \operatorname{Tr}_{H}\mathcal L_{p,q,n}(\rho,a)=0
 \quad\Longleftrightarrow\quad a=0.
\tag{4.6}\label{eq:joined-radical}
\]
In particular, at \(a=a_\rho\), the trace vanishes exactly on the critical
line.
\end{corollary}

This closes a frequent ambiguity.  Gamma normalization is not a missing
arithmetic source: it is a boundary coordinate on an already constructed
physical jet.  Nor is zero exclusion hidden in the decoder: off-line source
atoms are retained and produce nonzero unsquared trace.

\section{Reflected Green transport without atom loss}

The source also supplies a global reflected transport.  For a fixed prime
\(p\), let \(r_p(\rho)\) be the average of the exact adaptive prime-normal
reader at \(\rho\) and its functional-equation mate.  The arithmetic tail
estimate gives
\[
             |r_p(\rho)+1|\leq\frac14,
\tag{5.1}\label{eq:reader-quarter}
\]
therefore
\[
       \frac34\leq |r_p(\rho)|\leq\frac54
       \qquad\text{for every divisor copy }\rho.
\tag{5.2}\label{eq:reader-bounds}
\]
Multiplication by \(r_p\) defines an operator \(U_p\) on the complete
divisor Hilbert carrier.

\begin{theorem}[Reflected Green automorphism]
\label{thm:green-auto}
The operator \(U_p\) is boundedly invertible, with
\[
 \|U_p\|\leq\frac54,
 \qquad
 \|U_p^{-1}\|\leq\frac43,
 \qquad
 \|U_px\|\geq\frac34\|x\|.
\tag{5.3}\label{eq:green-bounds}
\]
It commutes with completed source Green reflection and with the physical
harmonic Hodge rotation.  If \(h\) is the fixed reflected Gaussian horizon,
then
\[
 (U_ph)(\rho)\neq0
 \qquad\text{at every divisor copy}.
\tag{5.4}\label{eq:green-atomic}
\]
The transported Green provenance vanishes globally if and only if the
Riemann hypothesis holds; locally, its coefficient at \(\rho\) vanishes if
and only if \(a_\rho=0\).
\end{theorem}

\begin{proof}
The pointwise bounds \eqref{eq:reader-bounds} imply
\eqref{eq:green-bounds} for the diagonal multiplier and its reciprocal.
Reflection invariance of \(r_p\) gives commutation with source Green
reflection; scalar diagonal multiplication gives commutation with harmonic
Hodge rotation.  Both factors in
\((U_ph)(\rho)=r_p(\rho)h(\rho)\) are nonzero, proving
\eqref{eq:green-atomic}.  The Green coefficient is a product of this nonzero
horizon with the oriented horizontal defect.  Its radical is therefore
exactly \(a_\rho=0\), and the global statement follows by quantifying over
the divisor support.
\end{proof}

Theorem~\ref{thm:green-auto} solves reflection compatibility and stability.
It deliberately does not solve selection.  The symmetrized reader is
reflection-even, so it transports the Green defect faithfully rather than
annihilating it.

\section{Boundary-defect conservation}

We isolate the algebraic law governing every cyclic-to-paracyclic descent.
Let a degree-one source carry a cyclic rotation \(\tau\) with
\(\tau^2=1\).  For an integer clock \(n\geq2\), define the normalized
source--dagger target rotation
\[
 R_{n,s}(x,y)=\bigl(y,\delta_n(s)x\bigr),
 \qquad
 \delta_n(s)=\exp\!\left((2\Rea s-1)\log n\right).
\tag{6.1}\label{eq:target-rotation}
\]
Then
\[
                   R_{n,s}^2=\delta_n(s)\,\mathrm{id}.
\tag{6.2}\label{eq:paracyclic-square}
\]
The scalar \(\delta_n(s)\) is the radial mate gain: it measures the mismatch
between source inverse and analytic Hilbert adjoint for the half-density
character.

For an arbitrary linear landing \(L:S\to V\), with source and target
rotations \(\tau\) and \(T\), define the one-step equivariance defect
\[
                         D=TL-L\tau.
\tag{6.3}\label{eq:defect}
\]

\begin{theorem}[Two-step boundary-defect conservation]
\label{thm:defect-conservation}
For every \(\tau,T,L\), one has
\[
                         TD+D\tau=T^2L-L\tau^2.
\tag{6.4}\label{eq:two-step-defect}
\]
Suppose \(\tau^2=1\) and \(T^2=\delta\,\mathrm{id}\).  Then:
\begin{enumerate}[label=\textup{(\roman*)}]
\item if \(D=0\) and \(\delta\neq1\), then \(L=0\);
\item if \(B\subseteq V\) is a \(T\)-stable linear subspace,
      \(D(S)\subseteq B\), and \(L(x_0)\notin B\) for one source atom,
      then \(\delta=1\).
\end{enumerate}
\end{theorem}

\begin{proof}
Expanding \eqref{eq:defect} gives
\[
 T(TL-L\tau)+(TL-L\tau)\tau=T^2L-L\tau^2.
\]
Under the square hypotheses this becomes
\[
                         (\delta-1)L=TD+D\tau.
\tag{6.5}\label{eq:monodromy-boundary}
\]
If \(D=0\) and \(\delta-1\neq0\), scalar cancellation gives \(L=0\).
For (ii), the right side of \eqref{eq:monodromy-boundary} lies in \(B\):
the first summand does by \(T\)-stability and the second directly by
\(D(S)\subseteq B\).  If \(\delta\neq1\), closure of \(B\) under scalar
multiplication would imply \(L(x_0)\in B\), a contradiction.
\end{proof}

\begin{corollary}[Retained source--dagger boundary class]
\label{cor:boundary-class-rigidity}
Let \(L_s:S\to\C^2\), and let \(B_s\subseteq\C^2\) be stable under
\(R_{n,s}\).  If
\[
 R_{n,s}L_s-L_s\tau\quad\hbox{takes values in }B_s
\tag{6.6}\label{eq:defect-boundary}
\]
and \(L_s(x_0)\notin B_s\) for one source atom, then
\(\Rea s=1/2\).  A single source-defined construction with these properties
at every nontrivial divisor point implies RH.
\end{corollary}

\begin{proof}
Theorem~\ref{thm:defect-conservation} gives \(\delta_n(s)=1\).  Since
\(n\geq2\), the real exponential in \eqref{eq:target-rotation} equals one
only when \(2\Rea s-1=0\).
\end{proof}

The distinction between vector non-erasure and class non-erasure is now
load-bearing.  Theorem~\ref{thm:gamma-atomic} proves that the physical
landing is nonzero in the target.  It does not prove that this vector is
nonzero modulo the arithmetic boundary required in
\eqref{eq:defect-boundary}.  Conversely, strict equivariance is no longer a
plausible missing construction: Theorem~\ref{thm:defect-conservation}(i)
annihilates every off-line fibre algebraically, before topology is imposed.

There is no freedom left in choosing the most favourable boundary.  Define
\[
 B_{\min}(\tau,T,L)=\ran D+\ran(TD).
\tag{6.7}\label{eq:min-boundary}
\]

\begin{theorem}[Maximal quotient law]
\label{thm:maximal-quotient}
Assume \(T^2=\delta\,\mathrm{id}\).  Then:
\begin{enumerate}[label=\textup{(\roman*)}]
\item \(B_{\min}\) is \(T\)-stable and contains \(D(S)\);
\item every \(T\)-stable subspace containing \(D(S)\) contains
      \(B_{\min}\);
\item a landed atom survives in some admissible boundary quotient if and
      only if it survives in \(V/B_{\min}\);
\item if \(\tau^2=1\) and \(\delta\neq1\), then
      \(L(S)\subseteq B_{\min}\); if \(L\) is surjective, then
      \(B_{\min}=V\).
\end{enumerate}
\end{theorem}

\begin{proof}
The first defect range lies in \(B_{\min}\) by definition.  Target rotation
sends it to the second range, while
\[
 T\bigl(TD(S)\bigr)=\delta D(S)\subseteq\ran D,
\]
proving (i).  If \(B\) is stable and contains \(D(S)\), it also contains
\(TD(S)\), proving (ii).  Part (iii) follows by choosing \(B_{\min}\) in one
direction and using (ii) in the other.  For (iv), if \(L(x)\notin B_{\min}\),
Theorem~\ref{thm:defect-conservation}(ii) would give \(\delta=1\).  Thus the
range is contained in \(B_{\min}\) when \(\delta\neq1\); surjectivity then
gives equality with \(V\).
\end{proof}

We now apply this universal law without adding a divisor-dependent gauge.
Let
\[
 c_\rho=\XiC^{(r_\rho)}(w_\rho)\neq0,
 \qquad
 L_\rho(z)=c_\rho z\quad(z\in\C^2),
\tag{6.8}\label{eq:physical-two-leg}
\]
and let \(\tau(x,y)=(y,x)\).  The coefficient in
\eqref{eq:physical-two-leg} is precisely the first nonzero physical
M\"obius--theta coefficient from \eqref{eq:gamma-landing}; hence the landing
is surjective and retains the exact multiplicity grade.  Take
\(T=R_{n,\rho}\) and form \(B_{\min,\rho}\) by
\eqref{eq:min-boundary}.

\begin{theorem}[Exact physical boundary classification]
\label{thm:physical-boundary-classification}
For every divisor point \(\rho\) and every \(n\geq2\),
\[
 B_{\min,\rho}=
 \begin{cases}
   0,&\Rea\rho=\tfrac12,\\
   \C^2,&\Rea\rho\neq\tfrac12.
 \end{cases}
\tag{6.9}\label{eq:boundary-dichotomy}
\]
If \(x_0=(1,0)\), then
\[
 [L_\rho(x_0)]\neq0\text{ in }\C^2/B_{\min,\rho}
 \quad\Longleftrightarrow\quad
 \Rea\rho=\frac12.
\tag{6.10}\label{eq:class-iff-line}
\]
Consequently, survival of these classes at every divisor point is equivalent
to RH.
\end{theorem}

\begin{proof}
If \(\Rea\rho=1/2\), then \(\delta_n(\rho)=1\), so
\(R_{n,\rho}=\tau\).  Scalar landing commutes with \(\tau\), hence \(D=0\)
and \(B_{\min,\rho}=0\).  Off the line, \(\delta_n(\rho)\neq1\); the landing
is surjective because \(c_\rho\neq0\), so
Theorem~\ref{thm:maximal-quotient}(iv) gives
\(B_{\min,\rho}=\C^2\).  Finally \(L_\rho(x_0)=(c_\rho,0)\neq0\), which
proves \eqref{eq:class-iff-line}; global quantification is exactly RH.
\end{proof}

Theorem~\ref{thm:physical-boundary-classification} is both construction and
obstruction.  It constructs the coherent maximal quotient permitted by
defect exactness, but proves that its desired non-erasure clause is not an
intermediate lemma: for the physical two-leg landing it is RH itself.  Any
independent source law must therefore do more than select a clever boundary;
it must explain why the actual arithmetic source cannot occupy the off-line
case of \eqref{eq:boundary-dichotomy}.

\section{Why regular completion cannot supply the boundary class}

The existing constructions sharply constrain the missing arrow.

\begin{theorem}[Regular-completion exclusion]
\label{thm:no-go}
The following statements hold for the constructed source and boundary
systems.
\begin{enumerate}[label=\textup{(\roman*)}]
\item After physical-jet normalization, every Hellinger--Gamma depth vertex
      has boundary charge one, while its Hilbert norm tends to zero.  Hence
      the graph closure contains \((0,1)\) and the boundary graph is not
      closed.
\item In the coefficient product topology, any closed linear relation that
      retains every leading finite prime block is maximally vertical:
      \[
          \mathcal R=\dom(\mathcal R)\times\C^{\N}.
      \tag{7.1}\label{eq:max-vertical}
      \]
      It transports no output information.
\item The old oriented Calder\'on--Witt boundary vanishes on the dense
      arithmetic cross-prime source span.  Every continuous extension to the
      completed packet is therefore the zero operator.
\item Reflection averaging commutes with Green and Hodge transport and is
      boundedly invertible, but it is even under reflection.  It preserves
      the off-line Green coefficient instead of forcing it to vanish.
\item A finite resolved packet can cancel only inside one full grade.
      Degree-one orientation, degree-two positive curvature, and higher
      fixed-carrier responses cannot cancel one another merely because they
      share the same radical.
\end{enumerate}
\end{theorem}

\begin{proof}
For (i), let \(v_n\) be the normalized depth vertex.  Its boundary is one
for every \(n\), while factorial Hellinger damping gives \(v_n\to0\) in the
completed carrier.  Thus \((v_n,1)\to(0,1)\), proving nonclosedness.

For (ii), finite-block surjectivity places every finite-support output in the
vertical multivalued part.  Finite-support sequences are dense in the
product topology, and closedness makes the multivalued part all of
\(\C^\N\).  A linear relation with full vertical part is its domain times the
entire output space.

For (iii), continuity and density determine an extension uniquely.  Since
the algebraic boundary is zero on the dense source span, its only continuous
extension is zero.  Part (iv) follows from
Theorem~\ref{thm:green-auto}.  Part (v) is
Theorem~\ref{thm:source-descent}(ii).
\end{proof}

Theorems~\ref{thm:defect-conservation} and
\ref{thm:maximal-quotient} strengthen this list.  Topology cannot rescue
strict equivariance, and boundary choice cannot rescue quotient survival.
The maximal quotient is canonical; off line it is zero for the physical
two-leg landing.  A viable successor must therefore prove that the physical
paracyclic defect is null-homologous in the literal cyclic arithmetic source.
It cannot enlarge the quotient by changing completion or boundary.

\section{Exact status of the arithmetic route}

The route can now be stated without mixing proved producers with conditional
consumers.

\begin{theorem}[Proved route envelope]
\label{thm:route-envelope}
There exists a target-free arithmetic source with all of the following
unconditional properties.
\begin{enumerate}[label=\textup{(\roman*)}]
\item It lands cofinally in every completed-zeta jet at every divisor point.
\item It retains the first nonzero multiplicity jet and the corresponding
      division-free top conormal.
\item Its leading coefficient lands nontrivially in the completed
      Gamma/co-Poisson incidence at every divisor stalk.
\item The landing has exact Hellinger coface recovery, and its unsquared
      Hodge trace has radical precisely \(a=0\).
\item A reflected adaptive source transport is boundedly invertible,
      commutes with Green reflection and harmonic Hodge, and is pointwise
      nonzero at every divisor copy.
\item Every ordinary closed, bounded, reflection-averaged, product-topology,
      or positive scalar completion covered by
      Theorem~\ref{thm:no-go} fails to force survival of the odd boundary
      class.
\item The canonical minimal defect boundary exists and gives the largest
      admissible quotient.  For the physical two-leg landing it is zero on
      the critical line and the whole target off it.
\item The genuine rational-adelic degree-one cyclic bar maps to the physical
      two-leg landing before divisor specialization; an ordinary-prime cycle
      has trace one, its cyclic residue detects the leading jet, and the
      physical residue defect is exactly monodromy minus one times that
      trace.
\end{enumerate}
\end{theorem}

\subsection{Literal cyclic comparison and prime-cycle residue}

Let \(A_{\rm CCM}\) be the minimal unitization of the algebraic rational
adeles-class-space crossed product and let
\[
 C_1(A_{\rm CCM})=A_{\rm CCM}\otimes A_{\rm CCM},
 \qquad \tau(a_0\otimes a_1)=a_1\otimes a_0.
\]
The zero-adele endpoint trace \(\varepsilon_0:A_{\rm CCM}\to\C\) induces
the genuine degree-one cyclic trace
\[
 \varepsilon_0^{(1)}(a_0\otimes a_1)=\varepsilon_0(a_0a_1),
 \qquad \varepsilon_0^{(1)}\tau=\varepsilon_0^{(1)}.
\tag{8.1}\label{eq:cyclic-trace}
\]
Let \(g\) be the product of the real Gaussian and the characteristic
function of the integral finite adeles.  Its zero-adele value is one.  For
an ordinary prime \(p\), let \(u_p\in C_1(A_{\rm CCM})\) be the tensor whose
crossed-product scaling labels are \(p\) and \(p^{-1}\), both with
coefficient \(g\).

\begin{theorem}[Ordinary-prime cyclic identity]
\label{thm:prime-cycle}
For every ordinary prime \(p\),
\[
                         \varepsilon_0^{(1)}(u_p)=1.
\tag{8.2}\label{eq:prime-cycle-one}
\]
\end{theorem}

\begin{proof}
The two scaling labels multiply to the identity label.  Twisted convolution
therefore places the product in the identity coefficient read by
\(\varepsilon_0\).  Zero is fixed by rational scaling, so zero-adele
evaluation is unchanged by the twist.  It is multiplicative on the Bruhat
tensor core, and each copy of \(g\) has value one.  Hence the product has
trace one.
\end{proof}

Define, without a divisor parameter,
\[
 C(x)=\bigl(\varepsilon_0^{(1)}(x),0\bigr),
 \qquad
 \operatorname{Res}(z_0,z_1)=z_0+z_1.
\tag{8.3}\label{eq:comparison-residue}
\]
Thus \(C(u_p)=x_0=(1,0)\).  For every jet grade \(k\), physical rung \(m\),
and spectral parameter \(w\), set
\[
 L_{k,m,w}(x)=
   \mathcal{FL}(\Psi_{m,k})(w)\,C(x),
 \qquad
 L_{k,w}^{\infty}(x)=\XiC^{(k)}(w)\,C(x).
\tag{8.4}\label{eq:global-cyclic-comparison}
\]
All objects in \eqref{eq:comparison-residue}--\eqref{eq:global-cyclic-comparison}
are defined before a zero is selected.

\begin{theorem}[Cyclic-to-paracyclic residue transgression]
\label{thm:residue-transgression}
The comparison \eqref{eq:global-cyclic-comparison} has the following
properties.
\begin{enumerate}[label=\textup{(\roman*)}]
\item For every \(k,K,\rho,p\),
\[
 L_{k,K+N,w_\rho}(u_p)\longrightarrow
 L_{k,w_\rho}^{\infty}(u_p).
\tag{8.5}\label{eq:comparison-limit}
\]
At \(k=r_\rho\), the right side is exactly
\(L_\rho(C(u_p))=L_\rho(x_0)\) from
\eqref{eq:physical-two-leg}.
\item The untwisted cyclic defect is residue-exact:
\[
 \operatorname{Res}\!\left(\tau C(x)-C(\tau x)\right)=0.
\tag{8.6}\label{eq:cyclic-residue-zero}
\]
\item Grades remain separate:
\[
 \operatorname{Res}(L_{k,w_\rho}^{\infty}(u_p))=0
       \quad(k<r_\rho),
 \qquad
 \operatorname{Res}(L_{r_\rho,w_\rho}^{\infty}(u_p))\neq0.
\tag{8.7}\label{eq:filtered-detection}
\]
No diagonalizability, unitary Jordan action, or simplicity of a zero is
used.
\item Let
\(D_{n,\rho}=R_{n,\rho}L_{r_\rho,w_\rho}^{\infty}
-L_{r_\rho,w_\rho}^{\infty}\tau\).  Then for every source chain \(x\),
\[
 \boxed{\quad
 \operatorname{Res}(D_{n,\rho}x)
  =\bigl(\delta_n(\rho)-1\bigr)c_\rho
     \varepsilon_0^{(1)}(x).
 \quad}
\tag{8.8}\label{eq:residue-transgression}
\]
In particular,
\(\operatorname{Res}(L_\rho(C(u_p)))=c_\rho\neq0\).
\end{enumerate}
\end{theorem}

\begin{proof}
Part (i) is the cofinal source-to-Xi limit in
Proposition~\ref{prop:source-landing}, applied at every independent grade,
together with \eqref{eq:prime-cycle-one}.  At the leading grade,
\(\XiC^{(r_\rho)}(w_\rho)=c_\rho\), giving the physical landing.
For (ii), cyclic trace invariance gives
\(\varepsilon_0^{(1)}(\tau x)=\varepsilon_0^{(1)}(x)\); coordinate sum
annihilates the resulting anti-diagonal vector.  Equation
\eqref{eq:filtered-detection} is exactly \eqref{eq:first-jet} multiplied by
the trace-one identity.  Finally, if
\(a=c_\rho\varepsilon_0^{(1)}(x)\), then the paracyclic rotation sends
\((a,0)\) to \((0,\delta_n(\rho)a)\).  Subtracting \((a,0)\) and summing
coordinates gives \((\delta_n(\rho)-1)a\), proving
\eqref{eq:residue-transgression}.
\end{proof}

\begin{theorem}[Exact residue boundary criterion]
\label{thm:residue-boundary-criterion}
For every \(n\geq2\) and divisor point \(\rho\),
\[
 \operatorname{Res}(B_{\min,\rho})=0
 \quad\Longleftrightarrow\quad
 \Rea\rho=\frac12.
\tag{8.9}\label{eq:residue-boundary-iff}
\]
Consequently,
\[
 \bigl[\operatorname{Res}(B_{\min,\rho})=0
       \text{ for every divisor point }\rho\bigr]
 \quad\Longleftrightarrow\quad \RH.
\tag{8.10}\label{eq:global-residue-iff-rh}
\]
More generally, any linear functional which annihilates
\(B_{\min,\rho}\) and is nonzero on \(L_\rho(x_0)\) forces
\(\Rea\rho=1/2\).
\end{theorem}

\begin{proof}
On the critical line, \(B_{\min,\rho}=0\), so every linear residue
annihilates it.  Off the line,
\(B_{\min,\rho}=\C^2\), while \(\operatorname{Res}(1,0)=1\); annihilation is
impossible.  This proves \eqref{eq:residue-boundary-iff}.  Quantification over
the divisor is RH.  The final assertion follows identically: off line the
minimal boundary is the whole target, so a functional vanishing on it cannot
detect any landed vector.
\end{proof}

The construction closes literal cyclic provenance, independent
ordinary-prime detection, and multiplicity-compatible physical landing.  It
also answers the annihilation demand exactly rather than postulating it.  The
last theorem is not ``find a source,'' ``find a residue,'' or ``retain
multiplicity.''  Those objects now exist.  It is the source-null-homology
statement that the physical analytic-dagger defect in
\eqref{eq:residue-transgression} is a literal cyclic boundary, or otherwise
has zero source residue, before divisor specialization.  By
\eqref{eq:prime-cycle-one} and \eqref{eq:residue-transgression}, that
statement forces \(\delta_n(\rho)=1\) at every divisor and is precisely the
RH-bearing selection law.  Calling it a routine comparison lemma would be
circular.

\section{Comparison with classical routes}

Table~\ref{tab:routes} compares logical inputs, not historical priority.

\begin{table}[p]
\centering
\footnotesize
\begin{tabularx}{\textwidth}{>{\raggedright\arraybackslash}p{0.17\textwidth}
                              >{\raggedright\arraybackslash}p{0.24\textwidth}
                              >{\raggedright\arraybackslash}X
                              >{\raggedright\arraybackslash}p{0.23\textwidth}}
\toprule
Program & Typical decisive object & What is already available here & Different unresolved gate \\
\midrule
Hilbert--P\'olya & A self-adjoint operator whose spectral parameters encode zeros & Complete divisor carrier, reflected Green transport, and maximal source--dagger quotient & No spectral realization or self-adjointness is assumed.  The quotient is already total off line; need independent arithmetic class survival, not an operator. \\
\addlinespace
Weil positivity & Positivity of an explicit-formula quadratic form & Positive Gamma/Green consumers with exact critical-line radical and an unsquared defect quotient & Positivity squares horizontal orientation.  Need source-forced nontriviality of the unsquared quotient class. \\
\addlinespace
Nyman--Beurling and Hedenmalm density & Closure or cyclicity of a function system & Dense arithmetic common tests, source-to-Xi landing, and divisor-faithful evaluation & Density controls representation.  The normalized Hodge boundary is nonclosable, so closure alone cannot impose same-fibre dagger. \\
\addlinespace
Sonin/co-Poisson & Fourier/Mellin duality and invariant subspaces & Exact Gamma/co-Poisson incidence, Green primitive, reflected transport, cyclic residue, and minimal mixed-defect boundary & Reflection duality relates mate fibres.  Need source-null-homology of the physical paracyclic defect. \\
\addlinespace
CCM/Meyer & Quotient trace, absorption spectrum, or cohomological boundary & Literal cyclic source, trace-one ordinary-prime cycle, physical two-leg comparison, and exact residue transgression & Need the analytic-dagger defect to be a literal cyclic boundary; the available trace proves what its residue would be. \\
\addlinespace
Deninger-type cohomology & Frobenius flow with a Hodge star or weight formalism & Half-density character, mate gain, top conormal, filtered cyclic comparison, and maximal quotient law & The remaining null-homology assertion kills the measured monodromy defect and is equivalent to critical-line placement. \\
\bottomrule
\end{tabularx}
\caption{Logical comparison of the arithmetic source--dagger route with classical programs.}
\label{tab:routes}
\end{table}

The present route is therefore not a disguised Hilbert--P\'olya program.
It does not search for an operator and then demand self-adjointness.  It
first constructs arithmetic source atoms and their singular conormal
landing.  The strict identification of source inverse with analytic dagger is
in fact impossible off line.  A Hilbert object may appear only after their
defect has descended to zero in a source-defined quotient.  If quotient
cyclicity and conormal survival are proved at the arithmetic boundary,
unitarity follows; if imposed at the Hilbert level, it is RH in operator
language.

The route also differs from a positivity criterion.  The sign of
\(a_\rho\) is present in the unsquared Hodge trace
\eqref{eq:joined-trace}.  Positive Gamma--Green curvature has the same
radical but degree two.  Resolved-grade conservation forbids using the
positive consumer as the missing oppositely signed degree-one term.

Finally, the route explains why repeated source reconstructions and density
theorems did not finish zero exclusion.  Those results solve attainment and
vector faithfulness.  Theorem~\ref{thm:maximal-quotient} also solves the
optimal boundary choice.  Theorem~\ref{thm:residue-transgression} now solves
the literal comparison and computes its homological obstruction.  The
remaining selection law must prove, from the literal source, that the
physical paracyclic defect is cyclically null-homologous.

\clearpage
\section{The falsifiable final law}

The remaining statement is narrower than ``construct a cohomology theory
for the zeta function.''  The preceding transgression theorem has already
passed tests for pre-divisor definition, physical comparison, independent
detection, and multiplicity.  A candidate final law must pass this audit.

\begin{enumerate}[leftmargin=2.2em]
\item \textbf{Literal boundary provenance.}  Exhibit a source chain whose
      cyclic boundary maps to the physical defect, or derive its zero residue
      from an independent arithmetic boundary identity.  Declaring the
      target defect exact is not a source proof.
\item \textbf{Fixed maximal quotient.}  The resulting detector must
      annihilate canonical \(B_{\min}\), not a smaller fitted subspace.
      Theorem~\ref{thm:residue-boundary-criterion} forbids an off-line
      shortcut.
\item \textbf{Prime-cycle compatibility.}  The boundary law must apply to
      \(u_p\) without erasing \(\varepsilon_0^{(1)}(u_p)=1\).  Otherwise it
      destroys the independent source detector it is meant to use.
\item \textbf{Filtered compatibility.}  The law acts grade by grade on
      \eqref{eq:filtered-detection}; it may not mix multiplicity jets into a
      unitary Jordan representation or assume simple zeros.
\item \textbf{Off-line falsifier.}  Any model that reproduces only an entire
      functional equation, reflection symmetry, positive Gamma factors, or
      hypertranscendence but admits off-line quartets must fail at a named
      source step.  Otherwise the proposed law is fitted target structure,
      not arithmetic selection.
\end{enumerate}

These tests locate useful future work.  Building another faithful divisor
carrier, comparison, positive form, residue, boundary, or
reflection-equivariant completion does not attack the final theorem.  All
are already constructed.  A useful advance must prove source-null-homology
of the measured physical defect in \eqref{eq:residue-transgression}.
Theorem~\ref{thm:residue-boundary-criterion} makes clear that this is the
RH-bearing theorem itself, not downstream infrastructure.

\section{Conclusion}

The arithmetic source route has crossed a sharper frontier than a list of
equivalent criteria suggests.  A target-free M\"obius--theta source reaches
every completed-zeta jet.  Its complete resolution retains arbitrary
multiplicity and lands in a nonzero top conormal.  The leading physical jet
enters the Gamma/co-Poisson incidence coefficient-faithfully at every divisor
stalk, and its unsquared Hodge trace detects horizontal displacement exactly.
A reflected adaptive transport is boundedly invertible, Green- and
Hodge-equivariant, and pointwise non-erasing.

Thus source, landing, conormal, Gamma incidence, reflected transport, and
atomic non-erasure are not open.  Neither are they an RH proof.  Source
descent shows why: each constructed regular transport faithfully carries the
off-line atom as well.  Reflection averaging removes odd orientation;
positive scalarization squares it; Hilbert closure makes the normalized
boundary nonclosable; product closure becomes maximally vertical; continuous
extension of the old boundary is zero.

The two-step defect law and maximal quotient classification change the
frontier.  Strict identification of source inverse with analytic dagger is
algebraically overdetermined off line.  Their minimal stable defect boundary
is canonical; for the physical leading-jet landing it is zero on the line and
all of \(\C^2\) off it.  Thus neither topology nor boundary choice can create
the desired class.

The literal comparison is no longer missing.  A canonical ordinary-prime
cycle has trace one, maps to the ordered physical atom, and is detected at
the first nonzero jet while all lower multiplicity grades vanish separately.
Its cyclic residue kills the untwisted comparison defect.  Analytic dagger
twisting changes that zero into the exact scalar
\((\delta_n(\rho)-1)c_\rho\).  This transgression formula isolates the sole
surviving claim: the physical paracyclic defect must be null-homologous in
the literal arithmetic source.  Global residue annihilation is equivalent to
RH, and no result in this paper assumes it.

\begin{thebibliography}{9}

\bibitem{Beurling1955}
A.~Beurling,
\newblock A closure problem related to the Riemann zeta-function,
\newblock \emph{Proc. Natl. Acad. Sci. USA} \textbf{41} (1955), 312--314.

\bibitem{Connes1999}
A.~Connes,
\newblock Trace formula in noncommutative geometry and the zeros of the
Riemann zeta function,
\newblock \emph{Selecta Math. (N.S.)} \textbf{5} (1999), 29--106.

\bibitem{Edwards1974}
H.~M.~Edwards,
\newblock \emph{Riemann's Zeta Function},
\newblock Academic Press, 1974.

\bibitem{HLS1997}
H.~Hedenmalm, P.~Lindqvist, and K.~Seip,
\newblock A Hilbert space of Dirichlet series and systems of dilated
functions in \(L^2(0,1)\),
\newblock \emph{Duke Math. J.} \textbf{86} (1997), 1--37.

\bibitem{Titchmarsh1986}
E.~C.~Titchmarsh,
\newblock \emph{The Theory of the Riemann Zeta-Function}, 2nd ed.,
revised by D.~R.~Heath-Brown,
\newblock Oxford University Press, 1986.

\end{thebibliography}

\end{document}
